\documentclass[12pt]{article}
\usepackage{amsmath}
\usepackage{amssymb}
\usepackage{geometry}
\geometry{margin=1in}

\title{GED Mathematics Worksheet- Rational Expressions and Functions}
%\subtitle{Rational Expressions and Functions}
\author{}
\date{}

\begin{document}

\maketitle

\section*{Instructions}
Solve all 40 problems. Show your work where applicable. Simplify all answers completely.

\section*{Part 1: Multiplying Rational Expressions (Questions 1-10)}

\begin{enumerate}

\item Multiply: $\frac{3x}{4} \cdot \frac{8}{9x}$

\item Multiply: $\frac{5a^2}{6b} \cdot \frac{12b}{10a}$

\item Multiply: $\frac{x+2}{x-1} \cdot \frac{x-1}{x+3}$

\item Multiply: $\frac{2x^2}{5y} \cdot \frac{15y^2}{8x}$

\item Multiply: $\frac{(x+4)(x-3)}{(x+1)} \cdot \frac{(x+5)}{(x-3)}$

\item Multiply: $\frac{x^2-4}{x^2+5x+6} \cdot \frac{x+3}{x-2}$

\item Multiply: $\frac{6m^2n}{7} \cdot \frac{14}{18mn^2}$

\item Multiply: $\frac{x^2-9}{x^2-1} \cdot \frac{x+1}{x+3}$

\item Multiply: $\frac{a^2-a-6}{a^2+4a+3} \cdot \frac{a+1}{a-3}$

\item Multiply: $\frac{2x^2-8x}{x^2-16} \cdot \frac{x+4}{2x}$

\end{enumerate}

\section*{Part 2: Dividing Rational Expressions (Questions 11-20)}

\begin{enumerate}
\setcounter{enumi}{10}

\item Divide: $\frac{5}{8} \div \frac{3}{4}$

\item Divide: $\frac{x^2}{y} \div \frac{x}{2y}$

\item Divide: $\frac{3a}{5b} \div \frac{9a^2}{10b^2}$

\item Divide: $\frac{m^2n}{6} \div \frac{mn^2}{12}$

\item Divide: $\frac{x+2}{x-3} \div \frac{x+2}{x+5}$

\item Divide: $\frac{x^2-4}{x^2+6x+9} \div \frac{x-2}{x+3}$

\item Divide: $\frac{x^2-8x}{x^2+8x+16} \div \frac{x-8}{x^2-2x-24}$

\item Divide: $\frac{2x^2+5x+3}{x^2-1} \div \frac{2x+3}{x-1}$

\item Divide: $\frac{a^2-9}{a^2-4} \div \frac{a+3}{a-2}$

\item Divide: $\frac{x^2-5x+6}{x^2+x-12} \div \frac{x-2}{x+4}$

\end{enumerate}

\section*{Part 3: Understanding Functions (Questions 21-30)}

\begin{enumerate}
\setcounter{enumi}{20}

\item Determine if the relation $\{(1, 2), (2, 4), (3, 6), (4, 8)\}$ is a function. Explain why or why not.

\item Determine if the relation $\{(1, 3), (2, 5), (1, 7), (3, 9)\}$ is a function. Explain why or why not.

\item For $f(x) = 3x - 2$, find $f(5)$.

\item For $g(x) = x^2 + 4$, find $g(-3)$.

\item For $h(x) = 2x^2 - 5x + 1$, find $h(2)$.

\item For $f(x) = \frac{x+3}{x-1}$, find $f(4)$.

\item For $g(x) = 5x - 7$, find $g(-2)$.

\item For $h(x) = x^3 - 2x + 4$, find $h(3)$.

\item For $f(x) = \sqrt{x + 5}$, find $f(4)$.

\item For $g(x) = |2x - 3|$, find $g(1)$.

\end{enumerate}

\section*{Part 4: Evaluating Functions with Expressions (Questions 31-35)}

\begin{enumerate}
\setcounter{enumi}{30}

\item For $g(x) = 3x + 5$, find $g(4n)$.

\item For $h(x) = 2x^2 - 3x + 7$, find $h(4y^2)$.

\item For $f(x) = x^2 - 6x$, find $f(a+2)$.

\item For $g(x) = 5x - 1$, find $g(2m + 3)$.

\item For $h(x) = x^2 + 2x - 8$, find $h(x-1)$.

\end{enumerate}

\section*{Part 5: Functions and Graphs (Questions 36-40)}

\begin{enumerate}
\setcounter{enumi}{35}

\item Use the vertical line test to determine if the graph of $y = x^2$ represents a function.

\item Use the vertical line test to determine if the graph of $x = y^2$ represents a function.

\item For the function $f(x) = 2x - 3$, find $f(2)$ using the equation. Then verify your answer by substituting into the function.

\item Describe the behavior of a function during the interval $-2 \leq x < 1$ if the function increases throughout this interval.

\item For a piecewise function, explain why different intervals might have different behaviors (increasing, decreasing, or constant).

\end{enumerate}

\newpage

\section*{Answer Key}

\subsection*{Part 1: Multiplying Rational Expressions}

\begin{enumerate}

\item $\frac{2}{3}$

\item $a$

\item $1$ (with restrictions: $x \neq 1, -3$)

\item $\frac{3xy}{2}$

\item $\frac{(x+4)(x+5)}{x+1}$ or $\frac{x^2+9x+20}{x+1}$

\item $\frac{x+2}{1}$ or $x+2$ (with restrictions: $x \neq -3, -2, 2$)

\item $\frac{4m}{3n}$

\item $\frac{x-3}{x-1}$ (with restrictions: $x \neq 1, -1, 3$)

\item $\frac{a-2}{1}$ or $a-2$ (with restrictions: $a \neq -3, -1, 3$)

\item $x-2$ (with restrictions: $x \neq 0, 4, -4$)

\end{enumerate}

\subsection*{Part 2: Dividing Rational Expressions}

\begin{enumerate}
\setcounter{enumi}{10}

\item $\frac{5}{6}$

\item $2y$

\item $\frac{2b}{3a}$

\item $2m$

\item $\frac{x+5}{x-3}$ (with restrictions: $x \neq -2, 3, -5$)

\item $\frac{x+3}{1}$ or $x+3$ (with restrictions: $x \neq -3, 2$)

\item $\frac{x(x-6)}{x+4}$ (with restrictions: $x \neq 0, -4, 8$)

\item $\frac{x+1}{1}$ or $x+1$ (with restrictions: $x \neq 1, -1$)

\item $\frac{a-3}{a-2}$ (with restrictions: $a \neq 3, -3, 2$)

\item $\frac{x-3}{1}$ or $x-3$ (with restrictions: $x \neq 2, -4$)

\end{enumerate}

\subsection*{Part 3: Understanding Functions}

\begin{enumerate}
\setcounter{enumi}{20}

\item Yes, it is a function. Each domain element (1, 2, 3, 4) is paired with exactly one range element.

\item No, it is not a function. The domain element 1 is paired with two different range elements (3 and 7).

\item $f(5) = 13$

\item $g(-3) = 13$

\item $h(2) = -1$

\item $f(4) = \frac{7}{3}$

\item $g(-2) = -17$

\item $h(3) = 25$

\item $f(4) = 3$

\item $g(1) = 1$

\end{enumerate}

\subsection*{Part 4: Evaluating Functions with Expressions}

\begin{enumerate}
\setcounter{enumi}{30}

\item $g(4n) = 12n + 5$

\item $h(4y^2) = 32y^4 - 12y^2 + 7$

\item $f(a+2) = a^2 - 2a - 8$

\item $g(2m+3) = 10m + 14$

\item $h(x-1) = x^2 - 4x - 5$

\end{enumerate}

\subsection*{Part 5: Functions and Graphs}

\begin{enumerate}
\setcounter{enumi}{35}

\item Yes, $y = x^2$ represents a function. Any vertical line intersects the parabola at most once, so it passes the vertical line test.

\item No, $x = y^2$ does not represent a function. A vertical line at $x = 4$ (for example) intersects the parabola at two points, so it fails the vertical line test.

\item Using the equation: $f(2) = 2(2) - 3 = 4 - 3 = 1$. Verification: Substituting $x = 2$ into $f(x) = 2x - 3$ gives $f(2) = 1$. 

\item During the interval $-2 \leq x < 1$, the function increases as $x$ increases. This means that as you move from left to right along the graph in this interval, the $y$-values get larger.

\item Different intervals can have different behaviors because a piecewise function is defined by different rules on different intervals. For example, one part might follow a linear rule (constant slope), another might follow a quadratic rule (changing slope), and another might be constant (zero slope).

\end{enumerate}

\end{document}